\documentclass[10pt]{article}
\usepackage[letterpaper,margin=0.9in]{geometry}
\usepackage{times}
\usepackage{hyperref}
\hypersetup{colorlinks=true,urlcolor=blue,linkcolor=blue}
\pagestyle{empty}
\setlength{\parindent}{0pt}
\setlength{\parskip}{0.25\baselineskip}
\usepackage{enumitem}
\setlist{nosep,leftmargin=1.5em}

\begin{document}

Oren E. Livne\\
Pine Birch Analytics\\
35 Kelinger Rd\\
Churchville, PA 18966-1033\\
Tel. 312-533-9130\\
\href{mailto:oren.livne@gmail.com}{oren.livne@gmail.com}

\vspace{0.3em}
\today

\vspace{0.3em}
Editor-in-Chief\\
INFORMS Journal on Computing

\vspace{0.3em}
Dear Editor,

I am pleased to submit my manuscript, ``A Near-Linear-Time Solver for Bilevel Congestion Network Design via
Directed Nonlinear Laplacian Flow,'' for consideration as a Research Article in INFORMS Journal on Computing.

\textbf{The problem.} Bilevel network design---computing tolls, capacity expansions, or routing incentives
that minimize total system travel time subject to a user-equilibrium response---has been understood at the
level of sensitivity analysis since Tobin and Friesz (1988): the design gradient reduces to one
linearized-equilibrium solve per Newton step, independent of the number of design variables. In principle this
makes gradient-based bilevel design exact and scalable. In practice it has not been, because that sensitivity
solve is itself a direct factorization of the equilibrium Jacobian, and direct factorization grows
superlinearly on the irregular, poorly-separable networks that selfish-routing and congestion-pricing
applications actually present---road networks with hub-like structure, communication and overlay networks,
peer-to-peer topologies. The result is a one-to-two-order-of-magnitude gap between how large a network we can
\emph{solve} for equilibrium and how large a network we can \emph{design} for, even though the underlying
mathematics says the two should cost about the same.

\textbf{The bridge.} This paper closes that gap by recognizing that the directed user equilibrium, written as a
source-form nonlinear Laplacian flow, has a Newton Jacobian that remains a symmetric weighted graph Laplacian
on the active subgraph---despite the one-sided, rectified congestion law that traffic and communication
routing actually obeys. That one structural fact is the bridge between the transportation application and the
algorithmic core: it means every solve the bilevel design loop needs---the forward equilibrium \emph{and} the
design adjoint---reduces to a short sequence of near-linear-time Laplacian solves, using the same class of
combinatorial multigrid and approximate-Cholesky solvers that have transformed large-scale spectral graph
algorithms over the last decade. A smoothing homotopy in the dead-zone width handles the one complication (the
rectified law's ill-conditioning and shifting active set) without sacrificing that near-linear cost.

\textbf{Why I believe this belongs in IJOC.} The contribution is squarely at the intersection IJOC serves: a
computational bottleneck in a classical OR problem (bilevel congestion pricing / network design), diagnosed
precisely as a linear-algebra scaling issue, and resolved with near-linear numerical methods, validated with
the rigor a computing-focused readership expects rather than asserted from theory alone:

\begin{itemize}
\item The forward-solve scaling is confirmed on a topology-diverse corpus of 96 real irregular graphs spanning
four orders of magnitude ($1.1\times10^4$ to $6.0\times10^6$ arcs) drawn from SuiteSparse/SNAP, LAW, DIMACS10,
and Newman, not just a single synthetic family.
\item The comparison against direct factorization is a fair fight, not a strawman: alongside the standard
unsymmetric UMFPACK/COLAMD baseline, I benchmark against CHOLMOD under a METIS nested-dissection
ordering---the theoretically appropriate direct solver for the symmetric positive-definite Jacobian this
problem produces---and show the near-linear solver still wins with a widening margin, and completes at network
sizes where even that stronger direct solver exhausts memory.
\item The framework extends to multicommodity logit stochastic user equilibrium, validated on the Sioux Falls
and Anaheim benchmark networks with their published origin--destination matrices, where destination-bundled
toll design reduces total system travel time by $7.7$--$13.1\%$.
\item Every experiment is reproducible: code, unit tests, and reproduction scripts are publicly available at
\url{https://github.com/orenlivne/dnlf}, and I have prepared the accompanying software/data archive per IJOC's
reproducibility requirements.
\end{itemize}

We are candid about scope: the advantage is specific to poorly-separable, non-planar topologies---on
near-planar road networks, where nested dissection is already near-optimal, the method does not win. I
believe this framing, together with the depth of the empirical validation, makes the paper a good fit for
IJOC's standards. Thank you for your consideration.

Sincerely,\\*
Oren E. Livne

\end{document}
